% fig_dfig_crowbar_cusum.tex — CUSUM crowbar detector response
\begin{tikzpicture}
\begin{axis}[
  thesis style,
  width  = 0.88\columnwidth,
  height = 4.8cm,
  xlabel = {Time (ms)},
  ylabel = {CUSUM statistic $g_k$ (pu)},
  xmin=0, xmax=120,
  ymin=0, ymax=5.2,
  xtick={0,20,40,60,80,100,120},
  ytick={0,1,2,3,4,5},
  legend pos=north west,
  grid=both,
  grid style={gray!20,line width=0.3pt},
]

%% Pre-crowbar CUSUM (slowly rising — SG-type DC transient)
\addplot[thick,thesisgray,dashed,domain=0:40,samples=60]
  {0.08*x};
\addlegendentry{Pre-crowbar (SG-path, slow rise)};

%% Post-crowbar CUSUM (flat, near zero — DFIG pure AC)
\addplot[thick,thesisblue,mark=*,mark size=1.5pt,mark options={thesisblue}]
  coordinates {
    (40,0.12)(50,0.08)(60,0.15)(70,0.06)(80,0.11)(90,0.09)(100,0.14)(110,0.07)(120,0.10)
  };
\addlegendentry{Post-crowbar (DFIG-path, $g_k \approx 0$)};

%% CUSUM threshold line
\addplot[thesisred,thick,dotted,domain=0:120,samples=2] {3.0};
\node[thesisred,font=\tiny,right] at (axis cs:62,3.12) {Threshold $h=3.0$};

%% Crowbar activation arrow
\draw[thesisorange,thick,->]
  (axis cs:40,2.0) -- (axis cs:40,3.2-0.15)
  node[above,font=\tiny,thesisorange,align=center]{Crowbar\\fired};

%% Crowbar event marker
\draw[thesisorange,thick,dashed,line width=0.8pt]
  (axis cs:40,0) -- (axis cs:40,5.2);

%% Trip decision annotation
\node[draw=thesisblue!60,fill=white,rounded corners=2pt,
      font=\tiny,align=center,inner sep=2pt]
  at (axis cs:90,1.8)
  {DFIG path active\\$\kappa_n = 2.1 < 30$\\trip confirmed};

\end{axis}
\end{tikzpicture}
